\section{Documented negative results}
\label{sec:appendix-negatives}
Six failed designs were tried and documented during the development of this
family (each with runnable or archived implementations and failure modes):
sparse-gather routing ($\sim$5$\times$ slower than dense masking on CPU);
a hierarchical (multi-scale pyramid) indexer (consistently worse than the
flat indexer); a multi-scale KV pyramid at 224 KV tokens
(overparameterized for the budget); a learnable block-bias offset (NaN
instability in $\sim$50\% of runs; root cause later traced to uninitialized
memory---`torch.empty' biases---and fixed); a linear-pool indexer with
gated residual KV (loses the per-head $\times$ per-dimension routing
flexibility of the softmax-overlap gate); and stacked $2{:}1 \rightarrow
2{:}1$ compression (harder to optimize within the budget). Negative results
of this kind are a deliberate part of the reproducibility package; they are
secondary to the paper's core claim.

\section{Surrogate attention (separate line of work)}
\label{sec:appendix-surrogate}
Distilling the attention block into a cheap inference-time student is a separate
line of work, deferred to a separate publication: it is \emph{not} part of this
paper's claims, and none of its numbers are apples-to-apples with the frozen-task
results in Section~\ref{sec:experiments} (they predate the frozen protocol). This
pointer exists only to acknowledge the overlap with attention-distillation work
\citep{goldstein2025radlads} and to record the retention decision: the surrogate
line, including the equal-budget small-transformer control and
distillation-epoch alignment, lives outside this manuscript.

\section{Reproducibility}
\label{sec:appendix-repro}
Every number in this paper traces to a repository artifact: the frozen task
specification (the sole authority on task definition), the pre-registered
factorial design, the run harness and its machine-readable results manifests
(125 factorial/baseline runs at $n{=}25$ + 180 cross-dataset runs + external
reference-point runs, JSONL), a verified
per-module parameter registry, causality-verification and wiring-audit
suites (all variants pass future-invariance and exact-zero-future-gradient
checks; the Normal baseline is bit-identical to a naive causal reference),
and a sanity-test suite covering shapes, gradients, parameter counts, and
tiny-task fitting (including a per-cell index-to-block mapping regression test).
Supplementary audits ship as machine-readable logs alongside the manifests: last-query
equivalence, the router gradient audit, canonical last-query latency with paired
deltas, the index-to-block mapping audit (with the pre-fix misrouting finding and
the re-run of the affected cells), and the unrounded 25-seed factorial
recomputation. Reported means and deltas are computed from the audited
manifests; the figure-generation scripts in this paper's source tree read
the same manifests, so every figure is regenerable. The freeze-clean
package --- environment lockfile, dataset SHA-256 manifests, automated
figure/table regeneration scripts, and runnable test suites (causality,
gradient masks, wiring audit, data integrity) --- is all in the repository.
The source tree is under git with a freeze baseline (commit recorded in
the project changelog); the remaining reproducibility steps
at release time are a tagged archive with a DOI (e.g., Zenodo) and the
official venue style swap.

\section{Future work}
\label{sec:appendix-roadmap}
Scope completed in this study versus open follow-ups, by axis:
\begin{enumerate}
  \item \textbf{Statistics.} Factorial-level inference is complete:
        $n{=}25$ paired seeds; paired Wilcoxon with
        Holm correction over cell-vs-Normal contrasts
        (Section~\ref{sec:stats}); per-seed main-effect contrast for H1.
        Deferred with justification: mixed-effects model (the $2\times2$ already
        identifies the three fixed effects; deferred as a robustness re-expression),
        HAC/block bootstrap (seed-level analyses), TOST/ROPE (no
        equivalence claims); temporal-dependence-aware statistics for the
        cross-dataset suite remain future work.
  \item \textbf{Datasets.} The suite covers
        the ETT family (ETTh1, ETTh2, ETTm1, ETTm2), large multivariate
        (Electricity, Traffic, Weather), different dynamics (Exchange,
        Solar-Energy), and three mechanistic synthetic families (sinusoids,
        change points, delayed-copy), all under the frozen protocol with 5
        paired seeds (180 runs; per-seed tables in the released results
        manifests). Remaining: horizon ($H \in \{192, 336, 720\}$)
        and lookback ($L \in \{512, 720, 1024\}$) sweeps.
  \item \textbf{Baselines.} DLinear \citep{zeng2023dlinear} and PatchTST \citep{nie2023patchtst}
        under the identical frozen protocol as external reference points,
        not peer methods (Section~\ref{sec:refpoints}): DLinear on all
        twelve datasets; PatchTST on the eight datasets where its
        channel-independent training fits the single-machine CPU budget
        (declared omissions on solar/electricity/traffic/weather).
        Remaining as future work: iTransformer, Informer, TimesNet, and an
        SSM baseline.
  \item \textbf{Hardware/latency.} Development-machine measurement is complete:
        Mean/p90/p99 forward latency with bootstrap
        CIs, throughput, analytic MACs, peak-memory accounting, training
        time, and batch/lookback scaling sweeps (batch 1--256,
        $L \in \{128, \dots, 4096\}$); see Section~\ref{sec:efficiency}.
        Remaining as future work: cross-hardware transfer (x86 CPU, GPU)
        and depth/width ablations --- conscious scope exclusions of this
        study.
  \item \textbf{Literature re-audit:} the 2026-08-16 sweep of the
        2024--2026 routing/compression-attention literature found no
        counterexample to the \emph{to our knowledge} claim
        (Section~\ref{sec:related}); a fresh sweep before locking the
        title remains good practice.
\end{enumerate}
